% Vectorruimten. Dürer, Melencolia I (1514): rechtsboven hangt een magisch vierkant van 4 bij 4. De magische
% vierkanten vormen een vectorruimte (het probleem van hoofdstuk 2 gaat over die van 3 bij 3).
% Motief: coördinaten ten opzichte van een basis. Twee basisvectoren u = (1,6; 0,45) en v = (0,55; 1,65) (in cm)
% vanuit O = (15,4; 3,6), het scheve rooster i u + j v voor i = 0..2 en j = 0..3, en de vector w = 2u + v,
% die je bereikt met twee stappen u en één stap v (gestippeld). De coördinaten van w zijn (2, 1).
% De bovenrand is zelf een vector: de lineaire combinatie 0,22 sin(t) + 0,14 sin(2t) van twee functies.
\def\omslagbestand{omslag-melencolia.jpg}
\def\omslagbron{Albrecht Dürer, \textit{Melencolia I}, 1514 (publiek domein)}
\def\omslagsamples{200}
\def\omslagrand{10.6 + 0.22*sin(\x*360/10.5) + 0.14*sin(\x*720/10.5)}
\def\omslagmotief{%
  \randlijn
  \begin{scope}[shift={(15.4,3.6)}]
    \foreach \i in {0,1,2}{\draw[white,line width=.6pt,opacity=.55] ({\i*1.6},{\i*.45}) -- ({\i*1.6+3*.55},{\i*.45+3*1.65});}
    \foreach \j in {0,1,2,3}{\draw[white,line width=.6pt,opacity=.55] ({\j*.55},{\j*1.65}) -- ({2*1.6+\j*.55},{2*.45+\j*1.65});}
    \foreach \i in {0,1,2}{\foreach \j in {0,1,2,3}{\fill[white,opacity=.8] ({\i*1.6+\j*.55},{\i*.45+\j*1.65}) circle (1.2pt);}}
    \draw[white,line width=1pt,dash pattern=on 3pt off 2pt] (1.6,.45) -- (3.2,.9) -- (3.75,2.55);
    \draw[omslaglijn=1.6pt,-{Stealth[length=3mm]}] (0,0) -- (1.6,.45);
    \draw[omslaglijn=1.6pt,-{Stealth[length=3mm]}] (0,0) -- (.55,1.65);
    \draw[fel,line width=2pt,-{Stealth[length=3.4mm]}] (0,0) -- (3.75,2.55);
    \fill[omslagstip] (0,0) circle (2.2pt);
    \node[text=white,font=\small,inner sep=1pt,anchor=north] at (.85,.1) {$\vec u$};
    \node[text=white,font=\small,inner sep=1pt,anchor=east] at (.22,1.0) {$\vec v$};
    \node[omslaglabel,anchor=south east] at (3.55,2.75) {$2\vec u + \vec v$};
  \end{scope}
}
